\documentclass[11pt]{article}
\usepackage[margin=1in]{geometry}
\usepackage{amsmath,amssymb,amsthm}
\usepackage{booktabs}
\usepackage{hyperref}

\title{Mathematical Formulation of DPO Training\\for Preference-Conditioned Elicitation}
\author{Patrick Cooper}
\date{\today}

\begin{document}
\maketitle

\section{User Preference Model}

A user's latent type is a parameter vector
$\theta = (\gamma, \alpha, \lambda)$ drawn from independent priors:
\begin{align}
  \gamma &\sim \mathrm{Beta}(a, b), \quad \gamma \in (0,1] \label{eq:gamma} \\
  \alpha &\sim \mathrm{LogNormal}(\mu, \sigma), \quad \alpha \ge 0 \label{eq:alpha} \\
  \lambda &\sim \mathrm{Uniform}(1, 3), \quad \lambda \ge 1 \label{eq:lambda}
\end{align}
where $\gamma$ is a discount factor encoding patience, $\alpha$ controls risk
aversion, and $\lambda$ controls loss aversion.

\subsection{Prospect-Theory Utility}

Given a wealth outcome $w$ and reference point $w_{\mathrm{ref}}$, the
value function follows Kahneman--Tversky prospect theory:
\begin{equation}\label{eq:prospect}
  u(w;\, \alpha, \lambda) =
  \begin{cases}
    (w - w_{\mathrm{ref}})^{\,1/(1+\alpha)}
      & \text{if } w \ge w_{\mathrm{ref}}, \\[4pt]
    -\lambda \cdot |w - w_{\mathrm{ref}}|^{\,1/(1+\alpha)}
      & \text{if } w < w_{\mathrm{ref}}.
  \end{cases}
\end{equation}
The exponent $1/(1+\alpha)$ produces concavity over gains and convexity over
losses. Increasing $\alpha$ increases curvature, i.e., strengthens risk
aversion. The multiplier $\lambda \ge 1$ makes losses hurt $\lambda$ times
more than equivalent gains.

\subsection{Single-Period Expected Utility}

An allocation $a \in \Delta^K$ (the probability simplex over $K$ channels)
induces a portfolio return distribution. Let $\mu, \sigma^2 \in \mathbb{R}^K$
denote the per-channel return means and variances conditioned on the current
regime. The portfolio moments are:
\begin{equation}\label{eq:portfolio}
  r_p = a^\top \mu, \qquad
  \sigma_p^2 = a^\top \mathrm{diag}(\sigma^2)\, a.
\end{equation}
Expected utility is estimated by Monte Carlo. Draw $N$ samples
$r_i \sim \mathcal{N}(r_p, \sigma_p^2)$, compute simulated wealth
$w_i = w_{\mathrm{current}} \cdot (1 + r_i)$, and average:
\begin{equation}\label{eq:eu-single}
  \mathrm{EU}(a;\, \theta)
  = \gamma^{\,T_{\mathrm{rem}}} \cdot
    \frac{1}{N}\sum_{i=1}^{N} u(w_i;\, \alpha, \lambda),
\end{equation}
where $T_{\mathrm{rem}}$ is the number of rounds remaining and
$\gamma^{T_{\mathrm{rem}}}$ discounts the entire payoff by the user's
patience.

\subsection{Multi-Period Expected Utility}

To make $\gamma$ identifiable from binary choices, diagnostic scenarios
targeting the discount factor use a compounding multi-period formulation.
Over $H$ periods, wealth compounds as
$w_i^{(t+1)} = w_i^{(t)} \cdot (1 + r_{i,t})$ with i.i.d.\ returns, and
per-round utility is discounted:
\begin{equation}\label{eq:eu-multi}
  \mathrm{EU}_{\mathrm{multi}}(a;\, \theta)
  = \frac{1}{N}\sum_{i=1}^{N} \sum_{t=0}^{H-1}
    \gamma^{\,t}\; u\!\left(w_i^{(t)};\, \alpha, \lambda\right).
\end{equation}
A patient user ($\gamma \approx 1$) weights later rounds nearly equally,
favoring low-variance allocations that compound reliably. An impatient user
($\gamma \ll 1$) effectively cares only about $t = 0$, so the discount
factor drops out of single-period comparisons but remains identifiable in the
multi-period setting.


\section{Synthetic User Choice Model}

When presented with two allocations $a_A, a_B$, a synthetic user with
parameters $\theta$ responds via softmax-rational choice:
\begin{equation}\label{eq:choice}
  P(\text{choose } A)
  = \sigma\!\left(\frac{\mathrm{EU}(a_A;\,\theta)
                       - \mathrm{EU}(a_B;\,\theta)}{\tau}\right),
\end{equation}
where $\sigma$ is the logistic sigmoid and $\tau > 0$ is a temperature
controlling noise (bounded rationality). This choice model provides the
supervised signal for both the analytical posterior and DPO pair construction.


\section{Reward Decomposition}

The system decomposes the reward for a recommendation action $a$ given
environment context $x$ and user type $\theta$ as:
\begin{equation}\label{eq:reward}
  R(a, x, \theta)
  = R_{\mathrm{quality}}(a, \mathrm{env})
  + \beta \cdot R_{\mathrm{alignment}}(a, \theta).
\end{equation}
$R_{\mathrm{quality}}$ is an objective, user-independent score (a
regime-adjusted Sharpe-like ratio). $R_{\mathrm{alignment}}$ measures how
well the recommendation matches the user's true optimal action.
The weight $\beta$ governs the quality--alignment tradeoff and is realized
through curriculum phasing: Phase~1 sets $\beta = 0$ (quality only);
Phase~2 introduces alignment through contrastive type-conditioned pairs.


\section{DPO Pair Construction}

Training data consists of triples $(x, y_w, y_l)$: a prompt $x$, a
preferred (``chosen'') completion $y_w$, and a dispreferred (``rejected'')
completion $y_l$. Both completions are allocation texts serialized as
channel-percentage lists. Pair construction differs across curriculum phases.

\subsection{Phase 1: Quality Floor Enforcement}

Given a game state $s$, generate $n$ candidate allocations by pooling:
optimal actions $a^*(\theta_j)$ for several randomly sampled $\theta_j$,
Dirichlet-random points $a \sim \mathrm{Dir}(\mathbf{1}_K)$, and
Gaussian perturbations of the optimals. Score each candidate by
$R_{\mathrm{quality}}(a)$ and check the quality floor
(diversification, dominance, bankruptcy probability).

The prompt $x$ is a natural-language rendering of $s$ with no user profile.
\begin{align}
  y_w &= \arg\max_{a \in \mathcal{C}} R_{\mathrm{quality}}(a),
    \label{eq:p1-chosen} \\
  y_l &=
    \begin{cases}
      \text{a floor-violating candidate} & \text{if one exists,} \\
      \arg\min_{a \in \mathcal{C}} R_{\mathrm{quality}}(a) & \text{otherwise.}
    \end{cases}
    \label{eq:p1-rejected}
\end{align}

\subsection{Phase 2: Type-Conditioned Alignment}

Given a game state $s$ and a sampled user type $\theta$, the prompt $x$
includes a natural-language user profile (time horizon, risk tolerance, loss
sensitivity with numerical parameters). Candidates again include
$a^*(\theta)$ for the target user plus contrasting and random allocations.

For the target user, score each candidate by expected utility:
\begin{equation}\label{eq:p2-chosen}
  y_w = \arg\max_{a \in \mathcal{C}}\; \mathrm{EU}(a;\, \theta).
\end{equation}

Sample a contrasting user $\theta'$ such that $|\gamma - \gamma'| \ge 0.2$
or $|\alpha - \alpha'| \ge 0.3$, then set:
\begin{equation}\label{eq:p2-rejected}
  y_l = a^*(\theta'),
\end{equation}
the allocation that is optimal for a different user type. The rejected
response is not a degenerate action---it is a high-quality action for the
wrong user. This forces the model to attend to the user profile in the
prompt, since the same allocation text switches between chosen and rejected
depending on $\theta$.


\section{DPO Loss}

With constructed pairs $(x, y_w, y_l)$, the model is trained via the
DPO loss from Rafailov et al.:
\begin{equation}\label{eq:dpo}
  \mathcal{L}_{\mathrm{DPO}}(\pi_\phi;\, \pi_{\mathrm{ref}})
  = -\mathbb{E}\!\left[
    \log\sigma\!\left(
      \beta_{\mathrm{dpo}} \cdot \left[
        \log\frac{\pi_\phi(y_w \mid x)}{\pi_{\mathrm{ref}}(y_w \mid x)}
        - \log\frac{\pi_\phi(y_l \mid x)}{\pi_{\mathrm{ref}}(y_l \mid x)}
      \right]
    \right)
  \right],
\end{equation}
where $\pi_\phi$ is the trainable policy (base model with LoRA adapters),
$\pi_{\mathrm{ref}}$ is the frozen reference policy (the unmodified base
model), and $\beta_{\mathrm{dpo}} = 0.1$ controls the strength of the KL
divergence penalty against the reference.

Intuitively, minimizing~\eqref{eq:dpo} increases the log-probability of
$y_w$ relative to $y_l$ under $\pi_\phi$, but only insofar as the shift from
$\pi_{\mathrm{ref}}$ is justified by the preference signal. The
$\beta_{\mathrm{dpo}}$ term prevents the policy from collapsing away from the
reference distribution.

\paragraph{Why this is ``conditional'' DPO.}
The loss function~\eqref{eq:dpo} is identical to standard DPO. What makes the
training type-conditioned is the structure of $x$. In Phase~2, the prompt
contains the user profile (e.g., ``Time horizon: Very long-term, discount
factor 0.95; Risk tolerance: Low, alpha 2.00''). Because $y_w$ is optimal for
the user described in $x$ and $y_l$ is optimal for a different user, the
gradient signal forces the model to learn the mapping from profile text to
appropriate allocation ordering. Standard DPO over unconditional prompts
would see contradictory preference signals across user types;
conditioning on the profile resolves the contradiction.


\section{Model Architecture}

The base model is Qwen2.5-1.5B-Instruct, loaded in 4-bit NF4 quantization
(QLoRA). Only low-rank adapters are trained:

\begin{table}[h]
\centering
\begin{tabular}{ll}
\toprule
Component & Setting \\
\midrule
Quantization & 4-bit NF4, double quantization, bfloat16 compute \\
LoRA rank / alpha & 16 / 32 \\
LoRA dropout & 0.1 \\
Target modules & \texttt{q\_proj}, \texttt{k\_proj}, \texttt{v\_proj}, \texttt{o\_proj} \\
Optimizer & Paged AdamW 8-bit \\
Batch size (effective) & $1 \times 8$ gradient accumulation \\
Epochs & 3 \\
DPO $\beta$ & 0.1 \\
\bottomrule
\end{tabular}
\caption{QLoRA and DPO training hyperparameters.}
\label{tab:arch}
\end{table}

The curriculum is sequential: Phase~1 trains on quality-only pairs, producing
a checkpoint. Phase~2 loads the Phase~1 checkpoint and continues training on
type-conditioned pairs, so quality floor compliance is established before
alignment is introduced.


\section{Deployment: LLM-in-the-Loop Elicitation}

After DPO training, the model drives the full elicitation cycle. At each
round $t$:

\begin{enumerate}
  \item The model receives the serialized environment state and interaction
    history, then generates a diagnostic query: two allocation options
    $(a_A, a_B)$.
  \item A synthetic user responds via~\eqref{eq:choice}.
  \item The choice is appended to the history.
  \item The model generates an interim recommendation $\hat{a}_t$
    conditioned on all choices observed so far.
\end{enumerate}

After $N$ rounds, the final recommendation $\hat{a}_N$ is evaluated by
Spearman rank correlation against $a^*(\theta_{\mathrm{true}})$ (alignment
score) and by quality floor constraint checks (violation rate).

The model maintains no explicit posterior over $\theta$; preference tracking
is implicit in the in-context conditioning on the choice history. This is
compared against an analytical baseline that maintains an explicit particle
filter posterior updated via the softmax-rational likelihood
of~\eqref{eq:choice}, with queries selected by expected information gain:
\begin{equation}\label{eq:eig}
  q^* = \arg\max_q \;
  H[\theta \mid h_t]
  - \mathbb{E}_{r \sim p(r|q, h_t)}\!\left[
    H[\theta \mid h_t, q, r]
  \right].
\end{equation}

The analytical baseline provides an upper bound on elicitation efficiency;
the DPO study measures how much of this bound the LLM can recover.

\end{document}
